\documentclass[a5paper,10pt]{article}

% https://github.com/InDevRus/filippov-solutions
% Нумерация страниц
\pagenumbering{gobble}

% Настройка полей
\usepackage{geometry}
\geometry{left=1cm}
\geometry{right=1cm}
\geometry{top=1cm}
\geometry{bottom=1.5cm}

% Вставка таблицы
\usepackage{array}
\usepackage{tabularx}

% Инструменты выравнивания формул
\usepackage{amsmath}
\usepackage{mathtools}

% Диакритика в математике
\usepackage{amsfonts}

% Вычисления размеров на ходу
\usepackage{calc}

% Удобные записи производных
\usepackage[italic=true]{derivative}

% Настройки сносок
\usepackage{enotez}

% Длинные знаки векторов
\usepackage{anyfontsize}
\usepackage[b]{esvect}

% Вставка картинок
\usepackage{graphicx}

% Вставка красной строки
\usepackage{indentfirst}
\setlength{\parindent}{1em}

% Условия в командах
\usepackage{xifthen}

% Луа-скрипты
\usepackage{luacode}

% Настройка команд отдельно в математическом и текстовом режимах
\usepackage{mathcommand}

% Для повторения LaTeX кода
\usepackage{multido}

% Конфигурация отступов формул
\delimitershortfall-1sp
\usepackage{mleftright}
\mleftright

% Растягивание объектов
\usepackage{relsize}

% Обтекание текстом
\usepackage{wrapfig}
\usepackage{subcaption}
\usepackage{floatflt}

% Поддержка ввода кириллицы
\usepackage[english, russian]{babel}

% Настройка корневой позиции для компилятора
\usepackage{import}
\providecommand*{\ProjectRootPath}{..}

\import{\ProjectRootPath/preambles/macros/}{theorems}

\import{\ProjectRootPath/preambles/fonts}{font_setup}

% Знак градуса
\usepackage{siunitx}

\import{\ProjectRootPath/preambles/macros/}{operators}
\import{\ProjectRootPath/preambles/macros/}{redefinitions}

\import{\ProjectRootPath/preambles/fonts}{letters_setup}
\import{\ProjectRootPath/preambles/fonts/adjustments}{custom_superscripts}

\import{\ProjectRootPath/preambles/custom}{formatting}
\import{\ProjectRootPath/preambles/custom}{frames}
\import{\ProjectRootPath/preambles/custom}{list_setup}

% TODO: перевести команды на современный стиль объявления
\input{../preambles/plotting}

\begin{document}

$\br{x + 2y}y' = 1$.
Введём новую функцию $u\br{x} = x + 2y\br{x}$.
$y\br{x} = \dfrac {1} {2} u\br{x} - \dfrac {1} {2} x$.
$y'\br{x} = \dfrac {1} {2} u'\br{x} - \dfrac {1} {2}$.
В терминах $u = u\br{x}$ получаем новое уравнение 
$u \br{\dfrac {1} {2} u' - \dfrac {1} {2}} = 1$.
Функция $u\br{x} = 0$ не является решением, так что
$u' - 1 = \dfrac {2} {u}$.
$u' = \dfrac {2 + u} {u}$.
$\dfrac {u} {u + 2} u' = 1$,
причём функция $u\br{x} = -2$ является решением.

$\tint \dfrac {u\br{x}} {u\br{x} + 2} u'\br{x} dx = \tint \br{1 - \dfrac {2} {u + 2}} du 
= u - 2 \ln\vbr{u + 2} + C$.

$u - 2 \ln\vbr{u + 2} + C = x$.
$y + C = \ln\br{x + 2y + 2}$.
$x + 2y + 2 = C\tpow {e}{y}$.
Выразить $y\br{x}$ в элементарных функциях для произвольного $C$ не получится, так что общее решение выразим в обратной форме:
$x\br{y} = C\tpow{e}{y} - 2y - 2$.
В эти решения также входит $u\br{x} = \linebreak = x + 2y = -2$ при $C = 0$.

Частное решение, проходящее через точку $\tsys{& x = 0 \\& y = -1}$, задаёт $0 = C \tpow {e} {-1} - 2 \cdot \br{-1} - 2$, откуда $C = 0$.
Искомое частное решение равно $x\br{y} = -2y - 2$.

\begin{figure}[!h]
    \centering
    \begin{tikzpicture}
        \pgfplotsset{
            Graph/.style = {
                samples = 200,
                restrict x to domain = -5.5 : 5.5,
                restrict y to domain = -5.5 : 5.5,
            },
        }
        \begin{axis}[
            width = 12cm,
            ymin = -4.2,
            ymax = 4.2,
            xmin = -5.2,
            xmax = 5.2, 
            xtick = {-10, ..., 10},
            ytick = {-10, ..., 10},
            minor tick num = 3,
            declare function = {
                f(\C, \x) = \C * exp(x) - 2 * x - 2;
            }]

            \foreach \C in {-2, -1.5, ..., 2, 0.2, 0.35, 0.7, 3, 5, 9, -5, -10, -20, -40, -80, -160} {
                \addplot[Graph, domain = -4.5 : 4.5]
                ({f(\C, x)}, {x});
            }
        \end{axis}
    \end{tikzpicture}
\end{figure}

\end{document}
